\section{Conclusion}
\label{sec:conclusion}

This work shows that query expansion in hybrid retrieval should account for
the distinct matching behavior of dense and sparse retrieval channels. DESA implements
this channel-asymmetric view using a shared set of generated reference
passages. Orthogonal residual expansion adds complementary semantic directions
to the dense query, whereas score-product anchoring uses the same generated
passages to reorder documents within the original sparse support. To evaluate this design
without tying its measured effect to a fixed top-$L$ cutoff, we determine the
fused ranking from complete channel lists and record the per-channel replay
stopping depths needed to certify its ordered top-$K$.

Across seven BEIR datasets, DESA improves equal-dataset macro nDCG@10 and
Recall@20 over Original by 3.82\% and 2.38\%, while reducing dense and sparse
replay stopping depths by 36.90\% and 36.56\%. The operator analysis clarifies the source
of these gains: sparse anchoring provides the more consistent effectiveness
improvement, and dense residual expansion contributes additional gains while
reducing the ranked evidence required for fusion. Under matched generated
references, paired tests do not detect a statistically significant
effectiveness difference between DESA and QuDAR-simple RRF. Comparing their
equal-dataset mean entry counts, DESA uses 64.16\% fewer available fusion-side
rank entries.

Together, these findings point to a division of labor for query expansion in
hybrid retrieval. Generated reference passages broaden semantic coverage in dense
retrieval and, when anchored to the original query, improve ordering in sparse
retrieval without broadening lexical support. This division of labor captures
the central design principle of DESA: Dense expands, while Sparse anchors.
